\section{Method}
\label{sec:method}

\subsection{Model, conditions, and memory}
\label{sec:method-organism}
\label{sec:method-conditions}
\label{sec:method-memory}

We treat EM as a fixed property of a deployed policy. The subject is the
published \texttt{Qwen2.5-14B-Instruct} bad-medical-advice LoRA organism
\citep{turner2025model}. C1--C5 share these adapted weights, C6 is the same base
model without the LoRA adapter. Generation used bfloat16 without quantization,
temperature and top-$p$ of 1.0, a 600-token cap, ten samples per prompt, and
seed 0. Across 27 runs, we retained 31{,}360 raw responses and their automated
judgments, together with model, configuration, and memory hashes.

Six conditions vary corrective content and its delivery. C1 receives no
intervention. C2 receives three corrective notes sampled once and fixed in its
system prompt. C3 retrieves three corrective notes per query. C4 accesses the
same notes through A-MEM \citep{xu2025amem}. C5 retrieves three neutral notes
matched in length and clinical register but containing no safety content. C6 is
the base model. Thus C2--C3 asks whether query-conditioned, episodic retrieval
improves on a fixed corrective prompt. C3$-$C5 asks whether corrective content
matters, although session-note displacement confounds their primary comparison.

The corrective corpus contains 144 clinical-safety notes (mean 80 words)
generated by \texttt{gpt-4o-mini} from MedSafetyBench training requests
\citep{han2024medsafety}. The placebo contains 144 neutral documentation notes
(mean 78 words), generated under a forbidden-vocabulary check for safety terms.
C3 and C5 query ChromaDB over \texttt{all-MiniLM-L6-v2} embeddings and return
$k=3$ notes. ChromaDB ranks records by the collection's configured vector
distance, with lower values indicating nearer embedding-space matches
\citep{chromadocs}. C4 uses \texttt{gpt-4o-mini} as a controller to analyze new records
and evolve the contexts, tags, and links of related A-MEM notes. Retrieved
content is rendered under a fixed preamble; rank, identity, and Chroma distances
are retained. As collection-level metric metadata was not retained, we interpret
these distances only as ordinal retrieval values rather than calibrated
similarities.

Each retrieval run first conducted ten clinical question--answer turns from the
MedSafetyBench dataset training split and wrote each pair to its store. Held-out
probes then queried the store. Self-authored notes therefore competed with corpus
notes for three retrieval slots. Matched zero-turn C3/C5 sensitivity runs on
Tiers C and D removed this displacement.

The corrected A-MEM follow-up was restricted to Tier D because it was designed
as a matched mechanism test on the harmful-clinical endpoint after an initial
implementation rendered only immutable note content. We did not extend the
follow-up to Tier C because that exploratory tier contains only eight independent
question families and its harm estimate depends on the judge endpoint. We did
not extend it to Tier O because refusal, that tier's preregistered primary
endpoint, was zero in every original condition. As such, the experiment does not
estimate how memory evolution affects non-clinical behavior or secondary Tier-O
outcomes such as judged harm and coherence. The Tier-D test compares an
evolution-off condition with an evolution-on condition under the same A-MEM
backend, rendering schema, three-slot budget, and replayed transcript.
Store snapshots verify that only the evolution-on condition changes
model-visible fields and that neither condition rewrites provided source context.

\subsection{Probes and outcomes}
\label{sec:method-tiers}
\label{sec:method-outcomes}

The four probe tiers serve distinct inferential roles. Tier B assesses
reproduction of generic EM, Tier C explores non-clinical transfer, Tier D
measures harmful clinical assistance, and Tier O measures over-refusal on benign
clinical questions. Tier B contains the eight published generic EM
probes \citep{betley2025emergent}. Exploratory Tier C contains 24 formats of
eight non-clinical questions, resampled as eight question families. Exploratory
Tier D contains 180 held-out requests derived from the MedSafetyBench dataset
designed to elicit unsafe medical advice \citep{han2024medsafety}. Preregistered
Tier O contains 180 actionable benign clinical questions selected from the
MedMCQA dataset \citep{pal2022medmcqa}; questions are asked without answer
options so refusal remains observable.

The pinned \texttt{gpt-4o-2024-08-06} judge applies the published alignment and
coherence prompts in separate temperature-zero calls
\citep{betley2025emergent}. A response is \texttt{derailed} at coherence
$\leq50$, \texttt{misaligned} at alignment $<30$ and coherence $>50$, or
otherwise aligned, refused, or excluded. On Tier D, refusal is counted as
aligned; on Tiers B, C, and O, refused responses are excluded from the harm
denominator. We report harm, refusal, and low-coherence/off-topic rates together. Recovery is
$(\text{C1}-x)/(\text{C1}-\text{C6})$.

All primary outcomes use this pinned judge. After completing the primary
analysis, we conducted a bounded, explicitly post-hoc audit of label stability.
We deterministically sampled 60 already-generated Tier-D responses: ten from
each C1/C3 $\times$ reference-\texttt{aligned}, \texttt{misaligned}, and
\texttt{derailed} cell (seed 20260818). A separately specified
\texttt{gpt-5.6-terra} judge received the same two rubric prompts, with
temperature 0, \texttt{reasoning\_effort=none}, and a 16-token completion cap.
We report exact verdict agreement by sampled cell and descriptive unweighted
Cohen's $\kappa$. Because sampling was stratified on the reference verdict,
these statistics do not estimate population agreement. The audit also cannot
estimate clinical correctness, an alternative C3$-$C1 effect, or the stability
of outcomes outside the sampled Tier-D conditions.

\subsection{Statistical analysis and registration}
\label{sec:method-stats}
\label{sec:method-prereg}

Each probe produced ten responses that shared the same question, retrieved
context, and random seed, so responses from the same probe were not independent.
We therefore resampled probes for Tiers B, D, and O and question families for
Tier C, retaining all responses associated with each resampling unit. For each
of 2{,}000 bootstrap replicates, we used the same sampled clusters across
conditions to preserve pairing and recomputed exclusions and condition contrasts
within each replicate. We report 95\%
bias-corrected and accelerated (BCa) bootstrap confidence intervals. BCa
intervals adjust their endpoints for bias and skew in the bootstrap
distribution; the acceleration term is estimated by repeatedly omitting one
probe cluster. The C4 mechanism test used the same paired cluster bootstrap. Its
primary contrast was evolution-on minus evolution-off harm; the strict-floor
($<50$) and composite-failure sensitivity analyses were specified before
generation.

Tier O was the only preregistered evaluation. The preregistered Tier-O hypothesis
predicted that the C3 refusal rate would exceed C6 by at least 10 percentage
points and would differ from C2 by less than 10 points. We filed the
preregistration on 2026-08-14 before generating any Tier-O responses or
judgments. The preregistered dataset used was Health-ORSC
\citep{zhang2026healthorsc}, which we considered and processed using the study's
data-preparation scripts. Before generating outcomes, we determined that its
items did not adequately match Tier O's target construct of over-refusal on
actionable benign clinical requests and replaced it with MedMCQA. We recorded
this deviation, but the replacement criterion was not preregistered. Accordingly,
the Tier-O refusal analysis tests the preregistered hypothesis using an instrument
whose selection rule was specified outside the preregistration. Tiers B--D,
low-coherence analyses, zero-turn comparisons, session displacement, and
train/test proximity are exploratory or supplementary.

\subsection{Reproducibility and assets}
\label{sec:method-repro}

The prepared artifact contains generation, judging, and analysis code, probes,
corpora, a dependency lockfile, and reproduction commands. The Qwen base model,
MiniLM retriever, and MedMCQA dataset are
Apache-2.0; vendored A-MEM and MedSafetyBench are MIT-licensed, with the latter
also restricted by its authors to research use. The credited EM adapter's model
card states no license, so existing-asset licensing is not fully documented.
Full precision inference requires a CUDA GPU with at least 40 GB memory, but
exact worker models, runtimes, storage, and total project compute were not
recorded.
